\documentclass[11pt]{article}
\usepackage[margin=1in]{geometry}
\usepackage{amsmath,amssymb,amsfonts,bm,mathtools}
\usepackage{graphicx,booktabs,hyperref,titlesec}
\usepackage{wsun_macros}

\title{\textbf{DoD Basic Research Proposal / White Paper Skeleton} \\
\large Structure-Preserving Multiscale Mechanics and Scientific Machine Learning}
\author{\textbf{Principal Investigator:} WaiChing Sun, Professor \\
Department of Civil Engineering and Engineering Mechanics, Columbia University \\
\texttt{wsun@columbia.edu}}
\date{}

\begin{document}
\maketitle

\section{Scientific Vision, Objective, and Heilmeier Framing}
% 1. What is the fundamental 6.1 basic research problem and why do current paradigms fail in extreme DoD-relevant regimes?
The overarching objective of this proposed research is to establish a mathematically rigorous, thermodynamically consistent multiscale modeling framework capable of predicting the nonlinear response and failure of heterogeneous multiphase materials under extreme loading conditions. Whereas classical phenomenological closures break down outside narrow calibration regimes and unconstrained machine learning surrogates violate frame indifference and lose ellipticity during localization, our approach formulates constitutive laws directly on invariant Riemannian manifolds coupled with Sobolev-regularized operator learning.

\section{State of the Art and Fundamental Knowledge Gaps}
% Contrast existing schools of thought and pinpoint the exact mathematical/physical bottleneck.

\section{Overarching Hypothesis and Technical Approach}

\noindent\textbf{Central Hypothesis:} \textit{By constraining neural functional approximations to Lie-algebraic invariant manifolds and enforcing second-law dissipation inequalities via Sobolev derivative supervision, multiscale constitutive closures can achieve provable thermodynamic stability and orders-of-magnitude acceleration in extreme-regime boundary-value problems.}

\subsection{Research Thrust 1: Geometric and Thermodynamic Foundation}
\noindent\textbf{Hypothesis 1:} ... \\
\textbf{Technical Approach:} ... \\
\textbf{Verification Metric \& Risk Mitigation:} ...

\subsection{Research Thrust 2: Structure-Preserving Neural Operators and Active Learning}
\noindent\textbf{Hypothesis 2:} ... \\
\textbf{Technical Approach:} ... \\
\textbf{Verification Metric \& Risk Mitigation:} ...

\subsection{Research Thrust 3: Multiscale Verification and Extreme-Regime Validation}
\noindent\textbf{Hypothesis 3:} ... \\
\textbf{Technical Approach:} ... \\
\textbf{Verification Metric \& Risk Mitigation:} ...

\section{DoD Relevance, Impact, and Deliverables Schedule}
% Articulate how fundamental 6.1 discoveries transition to Army (ARO) / Air Force (AFOSR) / Navy (ONR) capabilities.

\end{document}
